\documentclass[11pt,a4paper]{article}

\usepackage[margin=1in]{geometry}
\usepackage{amsmath,amssymb,amsthm}
\usepackage{booktabs}
\usepackage{longtable}
\usepackage{hyperref}
\usepackage{xcolor}
\usepackage{listings}

\hypersetup{colorlinks=true,linkcolor=blue!50!black,urlcolor=blue!50!black,
  citecolor=blue!50!black}

\lstset{
  basicstyle=\ttfamily\small,
  columns=fullflexible,
  breaklines=true,
  frame=single,
  framesep=4pt,
  xleftmargin=0pt,
  aboveskip=1em,
  belowskip=1em,
}

\providecommand{\ket}[1]{\lvert #1 \rangle}
\providecommand{\bra}[1]{\langle #1 \rvert}
\newcommand{\Sz}{S_0}
\newcommand{\So}{S_1}
\newcommand{\CC}{\mathbb{C}}
\newcommand{\FF}{\mathbb{F}}
\newcommand{\lean}[1]{\texttt{#1}}

\theoremstyle{definition}
\newtheorem{standard}{Rule}

\title{A Machine-Checked Certification for the Six-Qubit Moment Orbits}
\author{Formalization notes}
\date{}

\begin{document}
\maketitle

\begin{abstract}
Fifty-one no-go modules of this development each carry a Lean~4 proof, checked
by the Lean kernel alone: their axiom footprint is contained in
\lean{[propext, Classical.choice, Quot.sound]}, with no appeal to
\lean{native\_decide} or \lean{sorry} in any proof or tactic position. This note
describes what was certified, the standard the certification was held to, the
architecture of the development, the handful of algebraic mechanisms that
between them close all fifty-one cases, and how to reproduce the recorded checks.
It also records the complementary positive side: orbit~36 is certified rigid at
the manuscript's table representative, while each of the other ten completed
realized orbits has a source-level exact $A/B$ weight-enumerator image and
kernel-checked attainment. It
also reports, in Section~\ref{sec:errors}, the reasoning errors that were made
and caught during the work, because the way they were caught is the most
transferable part of the exercise.

\medskip
\noindent\textbf{Scope, stated precisely.} Each no-go orbit carries two theorems. The
first, \path{no_go_of_moments}, takes its Knill--Laflamme rows as hypotheses
and derives a contradiction. The second, \lean{no\_code}, assumes only that the
two logical states are supported on the residue-$0$ and residue-$\delta$ fibers,
are normalized, and satisfy \path{SS.HasDistance ![c, d] 3}; every
Knill--Laflamme row is then \emph{derived} through
\path{SixQubit/CodeBridge.lean}. So the development does establish, inside the
kernel, that no exact $((6,2,3))$ code is supported on any of these fifty-one
orbits. What it does \emph{not} yet establish is that every code must be
supported on one of them: the reduction of an arbitrary code to one of the
sixty-six rank-compatible moment orbits, and the transport of other angle
records to the chosen representatives, are not formalized. Section
\ref{sec:openitems} states this boundary.
\end{abstract}

\tableofcontents

%======================================================================
\section{What is certified}

\subsection{The statement}

Fix $q$ and an angle vector $a \in (\mathbb{Z}/q)^6$. The residue fibers are
\[
  \Sz = \{x \in \FF_2^6 : \textstyle\sum_i a_i x_i \equiv 0 \bmod q\},
  \qquad
  \So = \{x \in \FF_2^6 : \textstyle\sum_i a_i x_i \equiv \delta \bmod q\},
\]
and a putative code supported on this orbit is a pair of unit vectors
\[
  \ket{0_L} = \sum_{x \in \Sz} c_x \ket{x},
  \qquad
  \ket{1_L} = \sum_{y \in \So} d_y \ket{y}.
\]
Distance three means the Knill--Laflamme conditions hold for every Pauli $P$ of
weight at most two:
\[
  \bra{0_L} P \ket{1_L} = 0,
  \qquad
  \bra{0_L} P \ket{0_L} = \bra{1_L} P \ket{1_L}.
\]
For each of the fifty-one orbits the development proves that this system has
no solution. The Lean statement is a theorem named
\lean{no\_go\_of\_moments} in each orbit's module, of the shape
\begin{lstlisting}
theorem no_go_of_moments (c d : SS.BitString 6 -> C)
    (hm : MomentSystem (normSq (c y0)) ... (normSq (d z15)))
    (E_ : offSum c d ef_ w_ = 0) ...
    (D_ : diagBot c ef_ w_ = diagTop d ef_ w_) ...
    : False
\end{lstlisting}
Here \lean{offSum} is the off-diagonal Knill--Laflamme sum for a flip mask and a
$Z$-dressing. The functions \lean{diagBot} and \lean{diagTop} are the two sides
of a diagonal condition. Finally, \lean{MomentSystem} contains the pure-$Z$
diagonal conditions, the two normalizations, and nonnegativity of every
probability.

\subsection{The certification standard}
\label{sec:standard}

Five rules were fixed at the outset and held for all fifty-one orbits.

\begin{standard}
The transcription of the two fibers is certified by \lean{decide} over all
sixty-four bit strings, not asserted. Each module proves
\lean{S0\_eq\_fiber} and either \lean{S1\_eq\_fiber} or
\lean{S1\_residue}; these state that the listed support is exactly the required
residue fiber.
\end{standard}

\begin{standard}
Moment data are derived, not assumed. Every probability identity used in a proof
is a named theorem whose sole hypothesis is \lean{MomentSystem}. One invocation
of \path{linear_combination}, with coefficients from an exact rational reduction
against the moment matrix, closes it.
\end{standard}

\begin{standard}
Most transcription equations were generated by script from verified matrices.
For those generated modules, incidence lists, Walsh sign patterns, and certificate
coefficients are computed in Python from the fiber data and emitted as Lean. The
repository does not claim a complete byte-for-byte generator for every historical
source module; the committed statements are checked independently by the kernel,
the manifest gate, and the moment-transcription gate.
\end{standard}

\begin{standard}
\label{rule:axioms}
Only the standard axioms \lean{propext}, \lean{Classical.choice}, and
\lean{Quot.sound} occur. In particular,
\lean{Lean.ofReduceBool}, the signature of \lean{native\_decide}, never appears.
\end{standard}

\begin{standard}
The build is clean: zero errors and zero warnings, with no linter suppressed.
\end{standard}

Rule~\ref{rule:axioms} is the load-bearing one. \lean{\#print axioms} is
transitive, so a single clean line for \lean{no\_go\_of\_moments} certifies the
entire dependency cone beneath it.

\subsection{The result}

\begin{center}
\begin{tabular}{@{}p{0.48\linewidth}p{0.46\linewidth}@{}}
\toprule
No-go orbits certified & 51 (every no-go orbit) \\
Headline no-go theorems & 51 \\
Lines in the orbit modules & 41\,521 \\
Theorems in the orbit modules & 12\,521 \\
Orbit-36 rigidity chain & 16 modules; table-level closure \\
Ten explicit realization tables & 10 modules, 1\,606 lines \\
A/B enumerator layer & universal transform; ten exact images and attainment \\
Audit entries in \path{SixQubit/Audit.lean} & 592 \\
Audit entries with a non-clean footprint & 0 \\
Build errors & 0 \\
Build warnings & 0 \\
Occurrences of \lean{sorry} in tactic position & 0 \\
Occurrences of \lean{native\_decide} in tactic position & 0 \\
\bottomrule
\end{tabular}
\end{center}

The certified orbits are exactly the fifty-one the manuscript excludes:
$5$, $6$ and $47$--$52$ at projective order three; $12$--$15$, $17$--$20$, $24$,
$41$--$44$ and $46$ at order two; $25$--$32$ and $62$--$63$ at order four;
$33$--$35$ and $37$--$40$ at order five; and $53$--$61$ and $64$--$66$ at orders
six through nine. Orbit $66$ lives in \lean{SixQubit/Order9.lean} for historical
reasons and all others in \lean{SixQubit/Orbit<n>.lean}. The remaining fifteen of
the sixty-six rank-compatible orbits all target the realized orders $\{2,3,5\}$ and
are not excluded. Eleven of them ($7$, $8$, $9$, $10$, $11$, $16$, $21$, $22$,
$23$, $36$, and $45$) have explicit realized codes certified here; the four
remaining cases $1$--$4$ are outside the scope of this development.
A per-orbit table, generated from the sources, is in
\path{results/orbit-table.tsv}.

\subsection{Toolchain}

\begin{center}
\begin{tabular}{ll}
\toprule
Lean & \lean{leanprover/lean4:v4.27.0} \\
Mathlib & \lean{fc278a179e431ed10630af63fc05481680c0c7bd} \\
\bottomrule
\end{tabular}
\end{center}

The pins are recorded in \path{lean-toolchain} and \path{lake-manifest.json}.
Nothing outside Mathlib is used mathematically; the other manifest entries are
Mathlib's own dependencies.

%======================================================================
\section{Architecture}

The development is layered, and the layering is what made the harder orbits
tractable.

\subsection{The \lean{SS} layer}

\path{SS/} (1\,608 lines) holds the problem-independent infrastructure: bit
strings, Pauli strings, the enumeration of admissible weight-$\le 2$ Paulis, and
the predicates that decide which Knill--Laflamme conditions are automatic for a
given support. Nothing here is orbit-specific.

\subsection{Transcription modules}

Each orbit has a module containing:
\begin{itemize}
\item the two fibers as explicit \lean{def}s with a \lean{decide} proof that they
  are the residue fibers;
\item the \lean{MomentSystem} structure for that orbit, generated from the exact
  rational moment matrix;
\item the derived moment facts, each a theorem over \lean{MomentSystem};
\item \lean{offSum}, \lean{diagBot}, \lean{diagTop}, the incidence/flip/sign
  lemmas (all by \lean{decide}), and the Walsh expansion lemmas;
\item the orbit's Knill--Laflamme relations, each proved by
  \lean{linear\_combination} from expansions with script-computed coefficients;
\item the headline \lean{no\_go\_of\_moments}.
\end{itemize}

\subsection{Algebra modules}

Thirteen algebra modules carry the mathematics, separated from the bit strings
entirely.  The first four served the hard orbits of the initial batch:
\begin{center}
\begin{tabular}{llr}
\toprule
Module & Serves & Lines \\
\midrule
\lean{SixQubit/Block33.lean} & orbit 33 & 327 \\
\lean{SixQubit/Pair29.lean}  & orbit 29 & 636 \\
\lean{SixQubit/Table31.lean} & orbit 31 & 490 \\
\lean{SixQubit/Block26.lean} & orbit 26 & 415 \\
\bottomrule
\end{tabular}
\end{center}


\noindent Nine more were added later. Shared layers include \lean{Triangle},
\lean{Frame}, \lean{Spinor}, \lean{ParityFour}, and \lean{Hafnian}. The remaining
ones are \lean{Trace5}, \lean{Flip5}, \lean{Kernel5}, and \lean{K5}, the last for
the rebuilt orbit~30.

These modules mention no bit strings at all. They take complex numbers and the
relations among them and return \lean{False} or an algebraic conclusion. This has
a practical consequence worth stating plainly: \emph{the mathematical core of a
hard orbit can be written and compiled before any of the transcription
plumbing exists}, which is how the risky part gets de-risked first. It also means
the algebra can be audited on its own, without wading through several hundred
\lean{by decide} lemmas.

\subsection{The audit module}

\lean{SixQubit/Audit.lean} imports every orbit module and runs
\lean{\#print axioms} on each headline theorem and on selected intermediate
lemmas. Its output, captured in \path{results/axiom-audit.txt}, is 592 lines.
Every footprint is contained in the three standard axioms
\begin{lstlisting}
'SixQubit.<Module>.<theorem>' depends on axioms:
  [propext, Classical.choice, Quot.sound]
\end{lstlisting}

%======================================================================
\section{The mechanisms}

Thirty-five of the fifty-one orbits fall to one of a small ladder of
mechanisms, each of which is driven by a script that finds the certificate and
emits the Lean. The ladder, in increasing order of strength:

\begin{enumerate}
\item \textbf{$k$-edge Walsh.} A flip mask whose incidence set is small enough
  that the dressing sign matrix forces one incidence product to vanish, while the
  moment polytope forces both endpoints to have positive probability.
\item \textbf{Support collapse.} A coordinate forced to zero by the polytope
  removes incidences, which exposes new forced zeros; iterate.
\item \textbf{Propagation.} Two-term relations chain across masks.
\item \textbf{Quadratic closure.} The squared-modulus consequences of the
  two-term relations, reduced against the moment system, sometimes linearize and
  pin the probabilities outright.
\item \textbf{Positive row sum.} A shared lemma turning a row of the moment
  matrix into a contradiction with nonnegativity.
\item \textbf{Phase-ratio transitivity.} The quantities
  $\rho_x = \overline{c_x} d_{\hat x}$ are linked across masks; connectivity
  forces a common phase.
\item \textbf{Star hub and mirror permanent.} A hub vertex plus the vanishing of
  a $2 \times 2$ permanent (not determinant --- see below).
\item \textbf{Case analysis on the active set.} Driven by the polytope's face
  structure.
\item \textbf{Hub without a star.} A vanishing-pair tree replaces the hub.
\item \textbf{Diagonal Knill--Laflamme conditions.} The last rung; orbit 25 was
  the first orbit that could not be closed without them.
\end{enumerate}

A recurring trap deserves its own paragraph, because it caught this work three
times. When the fibers admit an involution, a four-term Walsh sum collapses to
\[
  (1 + \sigma_x \sigma_y)\, P(x,\hat y) + (1 + \sigma_z \sigma_w)\, P(z,\hat w),
\]
which is a genuine cross relation only when \emph{both} sign products are $+1$
and is vacuous otherwise. With the two pairs' relative sign flipped one obtains
the $2 \times 2$ \emph{determinant} rather than the \emph{permanent}, and the
determinant version is satisfied by any rank-one grid, so it yields no
contradiction at all. Any argument that reads a permanent where the data supply a
determinant is worthless, and the failure is silent.

%======================================================================
\section{The four hard orbits}

Four orbits needed genuinely new arguments. In each case the argument obtained
here is substantially shorter than the one in the manuscript, and in two cases it
avoids machinery the manuscript needs.

\subsection{Orbit 33: a block reflection}

$a = (0,1,1,3,3,3)$, $q = 5$, fibers $12 + 12$. The first qubit carries angle
zero, so both fibers are closed under flipping it and split into six blocks of
two; write $c_i, d_i \in \CC^2$.

The moment system alone pins every block norm to $R = (3,2,2,1,1,1)/10$. The
diagonal conditions for $F$ on the zero-angle qubit together with $F Z_j$ on a
second qubit give $C f = J C g$, where $C$ is the $6 \times 6$ tail-sign matrix
($\det C = 160$) and $J = \operatorname{diag}(1,-1,\dots,-1)$; and
\[
  C^{-1} J C = 2 R \mathbf{1}^T - I,
\]
which is equivalent to the six matrix identities
\[
  c_i c_i^\dagger + d_i d_i^\dagger = 2 R_i \rho .
\]
The cross conditions make blocks $1,2$ collinear with block $0$ on both fibers,
which forces $c_j^\dagger d_0 = 0$ and $c_0^\dagger d_j = 0$ for $j = 3,4,5$.
Orthogonality in $\CC^2$ then gives $\rho = I/2$ and simultaneously
$\rho = \tfrac{10}{3} c_0 c_0^\dagger$, so $c_0 c_0^\dagger = \tfrac{3}{20} I$,
whose off-diagonal entry is zero while
$\operatorname{normSq}(p_0)\operatorname{normSq}(q_0) = 9/400$.

Everything is division free; the only case analysis is \lean{mul\_eq\_zero}
against explicit nonzero rationals.

\subsection{Orbit 29: a \texorpdfstring{$K_4$}{K4} of orthogonalities}

$a = (1,1,1,1,2,2)$, $q = 4$, fibers $16 + 16$, and the two fibers are
\emph{not} complementary --- a fact that breaks the tooling used elsewhere and
required a separate relation extractor (\path{tools/core/relgen.py}).

Both fibers are closed under $X_5 X_6$, so each splits into eight pairs. Read a
pair as a vector in $\CC^2$; then every two-term relation is an orthogonality
$C_m^\dagger D_n = 0$, and the twelve that survive realise exactly the six edges
of $K_4$ on the four top pairs that carry them.

Two observations shortened the proof considerably relative to the manuscript's
fifteen-face audit.
\begin{itemize}
\item The uniform identity
  $P_{\{n,n'\}} = Q_n + Q_{n'} - \tfrac14$
  collapses the face table into a single lemma with one hop.
\item One never picks a nonzero vector. Instead the common bottom direction is
  \emph{constructed} from the anchor:
  $u := (-\overline{v_1}, \overline{v_0})$.
  It is automatically nonzero, automatically orthogonal to $v$, and
  $u \perp D_n$ is literally $\det[D_n, v] = 0$. Only two case splits survive.
\end{itemize}

\subsection{Orbit 31: a \texorpdfstring{$4 \times 4$}{4 by 4} table}

$a = (1,1,1,2,2,2)$, $q = 4$, fibers $16 + 16$ and complementary. The last three
angles are all $2$, so both fibers are stabilised by the Klein four-group on those
qubits and split into four cosets of four. The sixteen positions form a
$4 \times 4$ table with rows $A, B_1, B_2, B_3$ and columns $e_0, \dots, e_3$.

The moment normal form --- all of it verified exactly against the moment system
--- is
\[
  p_1(\bar v) = p_0(v) \ \text{at every position},
  \quad
  \sum_e \beta_{i,e} = \tfrac14,
  \quad
  \alpha_e = \sum_i \beta_{i,e} - \tfrac18 .
\]
The first of these lets every active position be written
$d_{\bar v} = \rho_v \overline{c_v}$ with $|\rho_v| = 1$, and twelve two-term
relations make $\rho$ constant along \emph{columns}.

Two facts remove the phase from the Lean side entirely, which is what made this
orbit finishable:
\begin{itemize}
\item $v_e = \overline{d_{\hat A_e}}$: the kernel vector is literally the
  conjugate of the top amplitude at the complement of $A_e$. No division, no
  phase object.
\item Two columns share a phase exactly when $v_e a_f = v_f a_e$, a
  division-free complex equation one can \lean{by\_cases} on.
\end{itemize}

Family one then says the four bottom columns are dependent with coefficients
$v$, so pairing against each column gives \emph{four} equations
$\sum_e v_e M_{ef} = 0$, where $M_{ef}$ is the Gram matrix of the $B$ columns.
Family three pins $M_{ef} = \overline{v_e} v_f$ whenever the phases differ; the
division-free form of that statement, and the fact that closed the last gap, is
the unconditional identity
\[
  \bigl(M_{ef} - \overline{v_e} v_f\bigr)
  \bigl(v_e \overline{v_f} - a_e \overline{a_f}\bigr) = 0 .
\]
Its second factor packages $a_e \ne 0$, $a_f \ne 0$ and $\rho_e \ne \rho_f$ into
one side condition. A column alone in its phase class then has bracket $3/8$ and
dies at once, and the five branches of the phase partition close in turn.

For comparison, the manuscript closes this orbit by a computer-assisted rational
classification of sixteen support types.

\subsection{Orbit 26: the cross-anchor contraction}

$a = (0,1,1,2,2,3)$, $q = 4$, fibers $16+16$. Blocks pair into four groups of
two; the group $\{6,7\}$ is the anchor and carries probability $1/8 + 1/8$.

Fix an anchor block $j$, set $X_i = C_j^\dagger D_i$, and let
$R_g = (X_{2g}, X_{2g+1})$. Contracting the outer identity with that anchor gives
$R_g J R_h^T = 0$ for the three pairs $g < h$, with $J$ the \emph{symmetric}
antidiagonal form; three nonzero pairwise $J$-orthogonal rows in $\CC^{1\times 2}$
lie on one coordinate axis.

The axis, however, is chosen \emph{per anchor}, and an earlier version of this
argument silently used both anchors at once. That is a real gap: an adversarial
review produced explicit witnesses in which $C_6 \not\parallel C_7$ and the
counting closes with slack. The repair is to contract the outer identity with
\emph{different} anchors in the two slots, on the $a$ index with $D_7$ and on the
$b$ index with $\overline{C_6}$, which yields
\[
  R_g^{(7)} J R_h^{(6)T} + R_h^{(7)} J R_g^{(6)T} = 0 ,
\]
reading $a_g b_h + a_h b_g = 0$ in the mixed configuration --- whose only
solution has one of the six entries zero. Two caveats, both checked rather than
assumed: the cross relation is \emph{vacuous} in the aligned branch, which is
closed separately by a moment and Cauchy--Schwarz argument; and the fact
\[
  P_0 + P_2 + P_4 = P_1 + P_3 + P_5 = \tfrac38
\]
is load-bearing --- an ablation shows that dropping it makes the whole fact set
feasible again, precisely at the aligned point
$P = (0,\tfrac14,0,\tfrac14,0,\tfrac14,\tfrac18,\tfrac18)$.

%======================================================================
\section{Orbit 36: rigidity rather than nonexistence}
\label{sec:orbit36}

Fifty-one of the sixty-six moment orbits are excluded. Orbit~36 is not: it is
realized, by an exact $((6,2,3))$ code whose projective diagonal group is
$C_5$. What is certified about it is therefore a different kind of statement,
and a stronger one than a no-go. From the same raw material every no-go module
consumes---support conditions, the moment system, and Knill--Laflamme rows of
Pauli weight at most two---the code is pinned uniquely.

\subsection{The statement}

\lean{SixQubit/Orbit36Certified.lean} carries the final table-level theorem
\[
  \begin{aligned}
  &\lean{table\_representative\_rigid}:\\
  &\quad \text{support in the two table fibers}
    \ \wedge\ \text{normalization}\\
  &\quad \wedge\ \lean{SS.HasDistance ![c,d] 3}
    \ \Longrightarrow\ \mathcal{R}.
  \end{aligned}
\]
for the representative $a=(1,1,1,2,2,4)$ and residues $(0,1)$ printed in the
manuscript.  The conclusion $\mathcal{R}$ has four parts:

\begin{enumerate}
\item \emph{the magnitudes are forced}: the twelve bottom probabilities are
      $1/10$ on six states and $1/15$ on the other six, and likewise for the
      top fiber---no state of either fiber may vanish;
\item \emph{the code is unique}:
      \lean{CodeEquiv.Equivalent} relates the pair $(c,d)$ to the explicit
      canonical pair $(c_{36}, d_{36})$;
\item \emph{the $A$ weight enumerator is determined}:
      $A = \bigl(1,\,0,\,\tfrac{21}{25},\,0,\,\tfrac{291}{25},\,
      \tfrac{384}{25},\,\tfrac{79}{25}\bigr)$.
\item \emph{the $B$ weight enumerator is determined}:
      $B = \bigl(1,\,0,\,\tfrac{21}{25},\,\tfrac{584}{25},\,
      \tfrac{183}{5},\,\tfrac{984}{25},\,\tfrac{671}{25}\bigr)$.
\end{enumerate}

The hypotheses matter as much as the conclusion. They are exactly what an exact
distance-three code supplies and nothing else: no amplitude is assumed nonzero,
no phase is assumed generic, no basis is assumed chosen. The support conditions
are containments, $\forall s \notin \Sz,\ c\,s = 0$, so zero amplitudes inside
the fibers are permitted throughout---which is what makes the full-support
conclusion in part~(1) a genuine consequence rather than a restatement of an
assumption.

\subsection{What the equivalence means}

Part~(2) is stated with respect to \lean{CodeEquiv.LocalMonomial}, a qubit
permutation together with a bit flip and a diagonal $(\alpha_i,\beta_i)$ per
qubit with $|\alpha_i| = |\beta_i| = 1$, composed with a unit scalar on each
logical state and an optional exchange of the two logical states. This is
\emph{smaller} than local-unitary equivalence, so proving uniqueness at this
level is the stronger claim; the manuscript states the result up to local
unitaries and physical-qubit permutations and notes the proof achieves the
smaller relation, which is the one formalized here.

\subsection{The chain}

Sixteen modules form the complete chain:

\begin{center}
\begin{tabular}{ll}
\toprule
\lean{K4Kernel} & the $K_4$ kernel used by the phase argument \\
\lean{Orbit36Code} & the canonical pair, its fibers, and
                     \lean{canonical\_is\_code} \\
\lean{Orbit36} & fibers, moment system, \lean{magnitudes\_forced},
                 \lean{full\_support} \\
\lean{Orbit36Phase} & $A_1 = 0$, the triangle sum, the weight-two table \\
\lean{Orbit36Enum} & the enumerator machinery through weight three \\
\lean{CodeEquiv} & local-monomial equivalence and distance invariance \\
\lean{Orbit36Class} & the classification up to that equivalence \\
\lean{Orbit36Enum4} & weights four to six \\
\lean{MacWilliams} & the subset MacWilliams identity at $n = 6$ \\
\lean{Orbit36Final} & the headline, assembled \\
\lean{Orbit36Vacuity} & the non-vacuity witness \\
\lean{CodeOrbit36} & derives the old moment/KL interface from
                     \lean{HasDistance} \\
\lean{BEnumerator} & the complete seven-row $A\mapsto B$ transform \\
\lean{CodeEquivEnum} & equivalence composition and fixed-map $A/B$ invariance \\
\lean{Orbit36Transport} & the table-to-symmetric gauge map in the kernel \\
\lean{Orbit36Certified} & the direct table-level rigidity, $A/B$, and attainment theorem \\
\bottomrule
\end{tabular}
\end{center}

\subsection{Non-vacuity}

A rigidity theorem with unsatisfiable hypotheses is worthless in a way a no-go
is not: a no-go with unsatisfiable hypotheses is merely uninformative, whereas
a rigidity statement would be asserting properties of nothing.
The older moment-level headline has an explicit canonical witness in
\lean{Orbit36Vacuity}.  More importantly, the theorem
\lean{table\_representative\_attained} in \lean{Orbit36Certified} constructs an
actual code on the original table fibers by inverse-transporting the canonical
pair.  It proves support, both normalizations, the genuine \lean{HasDistance}
predicate, and the displayed $A/B$ pair.  Thus the direct table-level theorem
is non-vacuous without appealing to an informal gauge identification.

\subsection{Orbit 36 at the level of codes}

Before the code-level bridge was added, the orbit-36 magnitude, uniqueness, and
enumerator results took the moment system and twenty-two Knill--Laflamme rows as
\emph{hypotheses}. The fifty-one no-go
modules never did: each \lean{CodeOrbit$n$.no\_code} derives them from the
single statement that $![c,d]$ is a normalized distance-three code supported on
the two fibers. Orbit 36 was the one module in the development without that
bridge.

The difference is not presentational. Before it, the rigidity theorem was a
statement about amplitude pairs satisfying certain moment conditions; the step
from ``is an exact code in orbit 36'' to ``satisfies those conditions'' was true
but unproved. Only with the bridge does the theorem say what the manuscript
says.

It was found by asking a mechanical question of the whole tree: for each of the
eleven realized orbits, which theorems exist, what are their hypotheses, and do
those hypotheses include \lean{SS.HasDistance}? Scanning for that turned up
three occurrences of \lean{HasDistance} across the orbit-36 chain, all three in
a \emph{conclusion} and none in a hypothesis.

\lean{tools/generators/gencode36.py} closes it by the route
\lean{gencode9.py} used for orbit 66, whose fiber structure is the same. Two
differences: orbit 36 names its top states $z_0,\ldots,z_{11}$ rather than
complements, and the generator checks $z_i=\overline{y_i}$ before writing
anything; and its moment structure carries nonnegativity fields that orbit 66's
does not, supplied by $\lVert\cdot\rVert^2\ge0$. The conclusion of
\lean{rigid\_code} is lifted verbatim out of \lean{orbit36\_rigidity} rather
than restated, so the wrapper cannot claim more or less than the theorem it
feeds.

That first bridge still worked in the symmetric gauge
$a^\star=(1,1,1,1,2,3)$ with residues $(0,4)$, whereas the manuscript's table
uses $a=(1,1,1,2,2,4)$ with residues $(0,1)$.  The remaining interface is now
closed in Lean rather than delegated to a Python orbit comparison.
\lean{Orbit36Transport.weight\_transport} proves on all 64 basis strings that
the explicit permutation and two flips shift the residue by one; support,
norm, and \lean{HasDistance} are then transported with the required logical
swap. \lean{CodeEquivEnum.Equivalent.trans} composes this witness with the
symmetric-to-canonical witness.  The same file proves the $A$ and $B$
enumerators invariant under exactly this fixed permutation/flip/swap map, and
\lean{SS.MacWilliams.bco\_transform\_six} supplies the full MacWilliams
transform.  The final theorem in \lean{Orbit36Certified} assembles these
ingredients.

\subsection{The gate that did not exist}

The four independent checks of Section~\ref{sec:reruns} all key off the
\lean{no\_go\_of\_moments} headline and obtain their module list from the
reviewed 51-module inventory. Orbit~36 has no such headline and is in no such
inventory, so every one of them silently skipped it: the single largest body of
Lean in this repository was the only part with no transcription gate at all.

\lean{tools/verification/rigidity.py} closes that. It checks the symmetric
gauge and, independently, the exact original table vector and its two twelve-state
fibers; it pins the residue-shift and logical-swap direction; and it checks the
same no-smuggling, Pauli-weight, moment-row, probability-order, and non-vacuity
obligations as before.  Finally it parses the complete table-level theorem and
requires exactly the two support, two normalization, and distance hypotheses,
together with the magnitude, full-support, direct-equivalence, and $A/B$
conclusions.  The current gate has thirty-four assertions.

Two details are worth recording, because both are the kind of thing that makes
a checker agree with a wrong module.

First, the orbit identity is not a string comparison. The Lean modules use the
angle record $a = (1,1,1,1,2,3)$ at $q = 5$ with fibers at residues $0$ and
$4$; \path{tools/manifest.json} records orbit~36 as $a = (1,1,1,2,2,4)$ with
$\delta = 1$. Those are different vectors, and that they describe the same
orbit is a fact to be proved, not a discrepancy to be waved away. The checker
reduces both to a canonical form under the group generated by local flips
($a_j \mapsto -a_j$, under which both fibers shift together so $\delta$
survives), qubit permutation, unit scaling ($a \mapsto ua$, $\delta \mapsto
u\delta$) and exchange of the two logical states ($\delta \mapsto -\delta$).
Both reduce to $\bigl((1,1,1,1,2,2),\ \delta = 1\bigr)$.

Second, the moment rows are re-derived here from the fibers rather than shared
with \lean{momentgate.py}. Two checkers built on one parser fail together; that
already happened once in this repository, when a declared-row counter and a
parsed-row counter shared a fragile prefix, missed the same row, and agreed.

The original mutation suite carried fourteen corruptions and initially caught
only two of them. The reason is instructive and is the same failure again in miniature:
an anchor such as \lean{(ho0 : offSum c d ex1 w0 = 0)} occurs five times in
\lean{Orbit36Final.lean}---once in the headline and again in the lemmas it is
assembled from---so a whole-file replacement mutated a lemma and left the
headline pristine, and twelve mutations were testing nothing while reporting
success. Mutations are now applied inside the headline region only, and a
mutation that fails to inject is reported as inert rather than as passing.  The
suite now retains those fourteen cases and adds twenty transport and table-closure
corruptions: all thirty-four are caught, with zero missed and zero inert.

\section{Ten realized exact-image orbits}
\label{sec:realizations}

Of the sixty-six moment orbits, fifty-one are excluded and fifteen are
realized. Eleven of the fifteen have a known exact weight-enumerator image;
orbit~36 is one of them and is treated in Section~\ref{sec:orbit36}. This
section covers the other ten---orbits 7, 8, 9, 10, 11, 16, 21, 22, 23
and~45. Each carries an explicit pair of words and a kernel proof that the pair
is an exact $((6,2,3))$ code. In addition, each now has a top-level theorem
quantified over arbitrary amplitudes on the two complete fibres, assuming only
support, normalization and genuine \lean{SS.HasDistance}; it proves the exact
$A/B$ image claimed by the source, and a second direction proves every claimed
parameter is attained.

\subsection{One layer, ten tables}

The orbit-36 code module carried its own ring of algebraic integers, its own
embedding into $\CC$, and its own transport from a decided certificate to
genuine Knill--Laflamme matrix elements. None of that depends on the amplitude
table, so it is collected in \lean{SixQubit/Realize.lean} and the ten modules
supply only data: two support lists, two amplitude tables, and a scale.

The ring is $\mathbb{Z}[\sqrt2,\sqrt3,i]$, carried as eight integers. It is
small enough to decide with and still large enough for every realization here,
because each table is recorded \emph{up to a single real scale}:

\begin{center}
\begin{tabular}{lll}
\toprule
words & scale & entries \\
\midrule
sixteen real terms & $1/4$ & $\pm1$ \\
eight real terms & $1/4$ & $\pm\sqrt2$ \\
eleven complex terms & $1/12$ &
  $2\sqrt3\,\zeta_6^{\,k}$, $2\sqrt6\,\zeta_6^{\,k}$ \\
twelve complex terms (orbit 36) & $1/(2\sqrt{30})$ &
  $2\sqrt3\,\zeta_6^{\,k}$, $2\sqrt2\,\zeta_6^{\,k}$ \\
\bottomrule
\end{tabular}
\end{center}

Pulling the irrational part out as a scale is what keeps the ring small.
Orbit~36's amplitudes are $1/\sqrt{10}$ and $1/\sqrt{15}$, which would need
$\sqrt5$; written as $\sqrt3/\sqrt{30}$ and $\sqrt2/\sqrt{30}$ they do not.
For the ten modules here the scales are rational, so the normalization
$\overline{s}\,s\,N=1$ is closed by \lean{norm\_num}.

The certificate the shared layer consumes is smaller than orbit~36's. Being a
code requires exactly three things of every Pauli of weight at most two: the
two diagonal elements agree and both off-diagonal elements vanish. Which
Paulis carry a nonzero diagonal scalar is extra information, so it is not part
of the shared realization certificate \lean{checkKL}; the later universal
enumerator stacks derive whatever orbit-specific information they need. The
scale is a free parameter of the transport theorem: it
cancels from the off-diagonal equations and enters the diagonal one
symmetrically.

\subsection{Two independent computations of the same certificate}

\lean{tools/generators/genreal.py} writes all ten modules. Two habits are
deliberate.

First, the amplitudes and the orbit records are \emph{parsed} from the source
tables, never transcribed. Hand transcription is where this project's errors
have actually come from, and the same script that parses the base words also
applies the qubit permutation, the translation and the optional exchange of the
two logical words that place each base into its orbit.

Second, the generator re-checks the whole certificate---all $4096$ mask pairs,
filtered by weight---in Python, in the same eight-integer arithmetic the Lean
uses, and refuses to write a table that does not verify. The kernel
\lean{decide} is then a second, independent confirmation rather than the only
one. Both computations agree on all ten.

\subsection{What the kernel cannot check}

Each module decides its own certificate, so the kernel establishes that the
exhibited pair is a code. It cannot establish that the pair belongs to the
orbit the module is named for: \path{tools/manifest.json} is not a Lean object.
That is exactly the gap that let four no-go modules target the wrong residue
before \lean{gate.py} existed.

\lean{tools/verification/realized.py} closes it, with twelve assertions per
module. The orbit identity is compared by canonical form under the group
generated by local flips, qubit permutation, unit scaling and exchange of the
two logical words, not by string equality---equivalent angle records look
different, and orbit~36 is the standing example. Both fibers are recomputed
from $(a,q)$, the two supports are checked disjoint, and the amplitude norms
are re-derived in the same ring. Its \lean{--selftest} injects six corruptions
and requires every one to be caught.

\subsection{A cross-check worth recording}

The exact weight enumerators of the ten exhibited codes were computed
independently, in exact arithmetic, and compared with the images the manuscript
claims:

\begin{center}
\begin{tabular}{lll}
\toprule
orbits & $(A_1,A_2,A_3)$ & manuscript \\
\midrule
8, 11, 23 & $(1,0,0)$ & $\mathcal E_1^A=(1,1,0,0,15,15,0)$ \\
7, 9, 10, 16, 21, 22, 45 & $(0,1,0)$ & $\mathcal E_2^A=(1,0,1,0,11,16,3)$ \\
\bottomrule
\end{tabular}
\end{center}

Every one agrees, and every enumerator sums to $32$. This is independent
evidence for both sides: the generated codes are the ones the manuscript
intends, and the manuscript's table is right about them.

\subsection{Ten exact A/B images}

For orbits 10, 11, 16, 22 and 23 the exact image follows from a common finite
squeeze.

A projective-order-two orbit with $q=4$ has a diagonal gate acting logically as
$\mathrm{diag}(1,-1)$, so its square is scalar on any code in the orbit. That
square is the $Z$-string on the coordinates where $a$ is odd, and its sign on a
support word $x$ is $(-1)^{A_a(x)}$, which is $+1$ because $A_a(x)$ is $0$ or
$2$ modulo four on the two fibers. So the string fixes the code---a condition
over the two supports, decided by the kernel.

A Pauli fixing both logical states has projector trace $2$, so it contributes
$4$ to the coefficient at its own weight, and coefficients are sums of squared
norms, so that is a lower bound. The shadow inequality is the matching upper
bound on the sum of the first three. Nothing is left free:
$\mathrm{Aco}_w=4$ and the other two vanish. The six-qubit MacWilliams normal
form then determines every remaining coefficient.

\lean{SixQubit/Pin.lean} carries the squeeze; each generated
\lean{Orbit$n$Enum.lean} supplies its mask and applies it. The conclusion is
about \emph{every} code supported in those fibers, not only the exhibited one,
and the realization module supplies a code attaining it, so the image is a
single point.  The $A$ vectors are
\[
  (1,1,0,0,15,15,0) \quad\text{for orbits 11 and 23},
  \qquad
  (1,0,1,0,11,16,3) \quad\text{for orbits 10, 16 and 22}.
\]
The corresponding $B$ vectors are
\[
  (1,1,0,30,45,33,18),
  \qquad
  (1,0,1,24,35,40,27).
\]
This second half is not a copied table: \lean{BEnumerator.lean} defines the
Frobenius-side coefficient and derives all seven MacWilliams rows by summing
the already proved subset duality over subsets of each cardinality.  Each
\lean{Orbit$n$Enum.lean} packages the result as
\path{enumerator_pair_unique_fiber} and proves the pair attained by its
explicit realized code.

\subsection{The five orbit-specific closures}

The remaining five require more than the scalar-mask squeeze. Orbits 7 and 8
share the affine image
\[
 A=(1,\alpha,1-\alpha,0,11+4\alpha,16-\alpha,3-3\alpha),
\]
\[
 B=(1,\alpha,1-\alpha,24+6\alpha,35+10\alpha,40-7\alpha,27-9\alpha),
\]
with $0\le\alpha\le1$ for orbit~7 and $1/2\le\alpha\le1$ for orbit~8.
Orbit~8 is closed by complete-support and Hafnian reductions. Orbit~7 uses the
covariant $K_5$ tangent/projector analysis, including the collinear and
zero-anchor boundary branches. Its structural conclusion is stated using a
literal finite-sum partial trace of $P/2$, and proves that some physical
two-qubit marginal has matrix rank two; no transported normal-form matrix is
substituted for that marginal. Explicit circle families attain every parameter
in both closed intervals.

Orbit~9 is rigid at the enumerator level:
\[
 A=(1,0,1,0,11,16,3),\qquad B=(1,0,1,24,35,40,27).
\]
The full-fibre active-support classification rules out the two-active branch,
produces a common equal-bit pair in the four-active branch, and also proves the
manuscript's structural clause: in the uniform gauge an actual weight-two
$Z_rZ_s$ fixes both logical words.

Orbit~21 has the same rigid $A/B$ pair, but its forced string has weight four.
The corrected proof expands projector traces into computational-basis
autocorrelations. The kernel checks the 512 possible
$(X\text{-mask},\text{Walsh frequency})$ pairs, proves the residual Walsh sum
vanishes for arbitrary amplitudes on the complete fibres, and combines it with
the exact Farkas identity
\[
 \mathrm{residual}=32(\mathrm{Aco}_2-4)-\mathrm{positiveRhs},
 \qquad \mathrm{positiveRhs}\ge0.
\]
The older conditional forced-vanishing route is retained only as historical
infrastructure and is absent from the final dependency chain.

For orbit~45 the result is the interval
\[
 A=(1,0,t,1-t,12-t,15+t,3),
 \qquad
 B=(1,0,t,22+2t,39-4t,42-2t,24+3t),
 \quad \frac23\le t\le1.
\]
Complete-fibre Knill--Laflamme rows give the block masses and phase relations;
division-free Gram lemmas cover zero-amplitude boundary cases. Two explicit
real-amplitude charts cover $[2/3,3/4]$ and $[3/4,1]$ and return genuine
normalized distance-three codes at the requested parameter. Here ``complete
fibre'' means support is contained in the whole fibre; it does not assert every
amplitude is nonzero.

\subsection{The shared shadow inequality}

The upper bound used across these exact-image arguments is the zeroth Rains
shadow inequality $A_1+A_2+A_3\le1$.

Its usual derivation builds the spin flip
$\Pi\mapsto Y^{\otimes n}\Pi^{\mathsf T}Y^{\otimes n}$ and quotes positivity of
$\operatorname{Tr}(\Pi\widetilde\Pi)$. There is a shorter route. Summing the
alternating Pauli sum over the $Z$-mask first collapses it by the character sum
already in \lean{SixQubit/MacWilliams.lean}, and what survives factors:
\[
  \sum_{P}(-1)^{\mathrm{wt}\,P}\bigl|\operatorname{Tr}(P\Pi)\bigr|^2
  \;=\;
  2^{n}\sum_{a,b}\Bigl|\sum_{u}(-1)^{|u|}\psi_a(u)\,\psi_b(\bar u)\Bigr|^{2},
\]
a sum of squared norms, so nonnegativity is immediate and no operator,
transpose or matrix positivity appears. The identity was checked numerically
before any Lean was written---on the orbit-36 code it gives $128/25$, exactly
what the six-qubit MacWilliams normal form predicts, and it holds on random
non-code pairs too, confirming that it is an algebraic identity rather than a
consequence of the code property. \lean{SixQubit/Shadow.lean} proves it. Both sides are
computed down to a common quantity: one direction is an interchange of
summation at depth four, the other does the $Z$-sum first, and the surviving
index is a complement. The statement
\lean{alt\_sum\_nonneg} holds for every $n$ and every number of
logical states; specializing to $n=6$ and grouping by weight gives
\lean{shadow\_bound}, namely $A_1+A_2+A_3\le1$ for an exact
$((6,2,d\ge3))$ code, in the unnormalized form
$\mathrm{Aco}_1+\mathrm{Aco}_2+\mathrm{Aco}_3\le4$ that the six-qubit
MacWilliams normal form already in the development consumes.

The orbit-specific lower bounds, phase/support reductions and constructive
families described above now supply the other half; the shadow theorem is no
longer being presented as completion of only a shared prefix.

\section{A tooling bug that gated three orbits}
\label{sec:bug}

The scan tooling under-generated the admissible Pauli set. For a weight-one flip
$e$ it returned only $\{0\} \cup \{e_j\}$, seven strings, where the correct set
is $\{0, e_i\} \times \{0, e_j\}$, twelve; it thereby dropped every $Y_i Z_j$. For
$e = 0$ it returned nothing at all, dropping all twenty-two pure-$Z$ conditions.
The correct total over all $|e| \le 2$ is
\[
  154 = 1 + 6 \cdot 3 + 15 \cdot 9 ,
\]
the number of weight-$\le 2$ Pauli operators on six qubits, which is the check
that should have been run at the start.

Three of the last four orbits were blocked by this. Orbit 33's argument needs
$F \otimes Z_j$ for $F$ ranging over $I, X, Y, Z$ on the zero-angle qubit, and
three of those four families were missing.

\emph{No earlier certification was invalidated by this under-generation bug.}
Those proofs used a
\emph{subset} of the true conditions, which makes them harder, never wrong. The
fix lives in \path{tools/core/kl.py}, which enumerates all sixty-four dressings
and filters on $|\operatorname{supp}(e) \cup \operatorname{supp}(z)| \le 2$
rather than being clever.

%======================================================================
\section{Errors made, and how they were caught}
\label{sec:errors}

Five reasoning errors were made during this work and all five were caught before
anything unsound entered the development. They share one shape: \emph{asserting
that a relation has a certain form without checking that form against the
originals}. Recording them, and the detection methods, is more useful than
pretending the path was straight.

\begin{enumerate}
\item \textbf{A wrong $\delta$.} Choosing the second fiber by size alone made the
  search accept easy certificates for the wrong mathematical objects. A later
  $d_{\mathrm{LP}}$ cross-check happened to disambiguate some orbits then still
  outstanding, but it is not the definition of $\delta$ and did not revisit
  modules already marked complete. Orbits $30$, $53$, $55$ and $65$ therefore
  used residues of additive order two instead of the orders their projective
  logical gates demand. A collaborator's audit found this; all four were rebuilt
  with the defining value $\delta=q/m$, and the corrected fibers are disjoint
  from the old ones. The guard that should have existed from the start is now
  \path{tools/manifest.json} plus \path{tools/verification/gate.py}, which checks
  the paper parameters, residue and both support fibers. The separate
  \path{tools/verification/momentgate.py} parses exact rational linear forms and
  checks every declared moment row against the probability order used by its
  headline theorem.
\item \textbf{The collapse trap, three times.} A four-term sum was read as a
  cross relation when the sign products made it vacuous, or as a permanent when
  the data gave a determinant. Fixed by writing an extractor
  (\path{tools/core/relations.py}) that computes every two-term relation genuinely
  in a mask's row space.
\item \textbf{Orbit 26: per-edge matrix relations that do not exist.} The
  weight-one masks carry eight incidences with rank only four and admit
  \emph{no} two-term relation; only group-level sums exist.
\item \textbf{Orbit 31: a family-three statement pushed too strong.} Claiming
  each term vanishes separately gives $G_{ef} t_{ef} = 0$, which is false. The
  bottom sign is $(-1)^{z \cdot x}$ and the $A$ row sits at complementary
  last-three-bits from the $B$ rows, so an odd-weight dressing flips the $A$-row
  sign relative to the others.
\item \textbf{Orbit 31: a covering count with a counterexample.} The claim that
  covering six column pairs with three rows forces one nonzero entry per row is
  false. For example,
  \[
    Z_1=\{0,1,2\},\qquad Z_2=\{0,1,3\},\qquad Z_3=\{2,3\}
  \]
  covers all six pairs with empty intersection. The branch does close, but by
  the full permanent condition, not by counting.
\end{enumerate}

The detection methods, in increasing order of usefulness:

\begin{enumerate}
\item \textbf{Row-space membership.} Reduced row echelon form against the mask's
  dressing sign matrix answers ``is this relation actually available?'' directly.
  This caught error 3.
\item \textbf{Replay against the originals.} Take a solution of the reduced
  system and evaluate the \emph{raw} Knill--Laflamme conditions on it. Exactly two
  failed for error 4, and both were weight-one dressings, which pointed straight
  at the sign. This is cheap and should be run on every reduction.
\item \textbf{Adversarial review.} A multi-agent pass whose agents were
  instructed to \emph{refute} rather than confirm produced explicit witnesses for
  error 3's consequence in orbit 26, including a configuration satisfying every
  fact the argument used.
\item \textbf{Writing the Lean.} Error 5 surfaced while formalising the branch:
  the combinatorial step had to be made concrete, and the counterexample fell
  out. This is the cheapest of the four and argues for a working practice:
  \emph{do not derive at length on paper and formalise in one batch; formalise as
  you go.}
\end{enumerate}

For the final two orbits the work was delegated to multi-agent workflows under
three non-negotiable instructions: verify by machine rather than assert, actually
compile rather than report compiling, and classify every hypothesis of the
headline theorem. That last one is the important one. A headline
\lean{False} theorem is only as good as its hypotheses, and a proof becomes a
shell the moment an algebraic fact that should have been derived is smuggled in as
an assumption --- a defect invisible from the outside. Both final headlines were
classified by hand after the fact:

\begin{center}
\begin{tabular}{lrrrrr}
\toprule
Orbit & Hypotheses & \lean{offSum = 0} & \lean{diagBot = diagTop} &
  \lean{MomentSystem} & Other \\
\midrule
31 & 46 & 33 & 12 & 1 & 0 \\
26 & 41 & 24 & 16 & 1 & 0 \\
\bottomrule
\end{tabular}
\end{center}

%======================================================================
\section{Reproduction}

\subsection{Checking the proofs}

From the repository root:
\begin{lstlisting}
lake build
\end{lstlisting}
should report zero errors and zero warnings. Then
\begin{lstlisting}
lake build SixQubit.Audit
\end{lstlisting}
prints the 592 axiom lines; every footprint must be contained in
\lean{[propext, Classical.choice, Quot.sound]}. A stored copy of the expected
output is \path{results/axiom-audit.txt}, so
\begin{lstlisting}
lake build SixQubit.Audit 2>&1 | grep "depends on axioms" \
  | sed 's/^info: [^ ]* //' | diff - results/axiom-audit.txt
\end{lstlisting}
should produce no output.

To confirm the absence of the forbidden constructs:
\begin{lstlisting}
grep -rn "sorry\|native_decide" SixQubit/ SS/
\end{lstlisting}
Every hit must be inside a docstring or comment; the exact comment count is not
an invariant. The transitive \lean{\#print axioms} audit is
the load-bearing check: an active escape hatch would add \lean{sorryAx} or
\lean{Lean.ofReduceBool}, neither of which appears.

\subsection{Regenerating the Lean}

Most transcription modules were emitted or repaired by scripts in
\path{tools/generators/}. Their inputs vary: some read \path{tools/orbits/}, some
read \path{tools/manifest.json} or checked certificate data, and some algebra
layers are hand-written. The provenance index in \path{tools/README.md} names the
major routes and explicitly marks the historical order-nine module, for which no
complete raw-module generator is committed. The scripts require Python with
\texttt{sympy}, \texttt{numpy} and \texttt{scipy}. For scripts using the shared
modules, set \texttt{PYTHONPATH} to include \path{tools/core} and
\path{tools/orbits}, or run from the repository root with
\begin{lstlisting}
python -c "import sys; sys.path[:0] = ['tools/core','tools/orbits']; ..."
\end{lstlisting}

\subsection{Re-running the independent checks}
\label{sec:reruns}

\path{tools/verification/} holds the scripts that answer questions the Lean
proofs do not: whether an orbit is excluded numerically at all, which relations
are genuinely available, and which branches of a case analysis are live.
Highlights:
\begin{center}
\begin{tabular}{@{}p{0.34\linewidth}p{0.62\linewidth}@{}}
\toprule
Script & What it answers \\
\midrule
\path{feas2.py}, \path{feasnew.py} & Is the full weight-$\le 2$ system numerically feasible or infeasible, including positive and negative controls? \\
\path{gate.py} & Do the paper parameters, residue and two supports match? \\
\path{momentgate.py} & Is every declared moment row the exact intended linear form? \\
\texttt{momentgate.py --selftest} & Does the non-exhaustive mutation suite fail closed? \\
\path{classify.py} & Does every algebraic headline hypothesis have an allowed form? \\
\path{admis.py} & Does every headline KL row have Pauli weight at most two? \\
\path{rigidity.py} & Does orbit~36 meet the original-table code-level rigidity contract? \\
\path{realized.py} & Do all ten explicit codes match the manifest fibres and exact arithmetic? \\
Orbit 45 diagnostics & Independent Gram and two-chart checks in \path{o45_gram.py} and \path{o45_attain_check.py} \\
\path{close31.py} & Orbit 31 against the raw conditions, phases free \\
\path{branch31.py} & Which phase partitions of orbit 31 are excluded \\
\path{collapse29.py} & Orbit 29's support collapse \\
\path{stepA\_*.py} & The four silent failure modes behind orbit 26's Step A \\
\path{t26\_*.py}, \path{x26\_*.py} & Orbit 26 facts, ablations, branch tree \\
\bottomrule
\end{tabular}
\end{center}

\path{feas2.py} includes certified orbits as controls: an orbit with a
kernel-checked no-go proof must come out excluded, and it does (orbit 25 gives a
best residual of $0.79$). Running it without the controls would not be a test.

\subsection{The workflow records}

\path{results/workflows/} holds the scripts of the multi-agent passes and their
verdicts, including the adversarial review that found the orbit 26 gap and the
audits of the final two assemblies. They are kept because they contain the
negative results --- which relations are \emph{not} available, which branches are
vacuous --- that are expensive to rediscover and easy to assume away.

%======================================================================
\section{Repository layout}

\begin{longtable}{@{}p{0.30\linewidth}p{0.64\linewidth}@{}}
\toprule
Path & Contents \\
\midrule
\endhead
\path{SixQubit/} & 51 no-go stacks, 11 realized-orbit certification stacks, shared algebra, and the audit \\
\path{SS/} & problem-independent bit-string and Pauli infrastructure \\
\path{formalization/} & this process document in source and PDF form \\
\path{tools/core/} & fibers, Pauli enumeration, moment matrices, relation extractors \\
\path{tools/orbits/} & per-orbit data and block structure \\
\path{tools/generators/} & the scripts that emit the Lean modules \\
\path{tools/verification/} & numerical and symbolic checks independent of Lean \\
\path{results/} & axiom audit, orbit table, workflow records \\
\bottomrule
\end{longtable}

%======================================================================
\section{Open items}
\label{sec:openitems}

Two items concern the formalization itself.

\begin{enumerate}
\item \textbf{The classification is not formalized.} Neither the reduction of an
  arbitrary code to one of the sixty-six rank-compatible moment orbits, nor the
  equivalence of other angle records to the chosen representatives, has a Lean
  proof.
\item \textbf{The exact images of orbits 1--4 are not formalized.} Their source
  proofs are still being completed. The present development therefore closes
  all 51 no-go orbits, orbit~36 at rigidity level, and the ten WE orbits
  $7,8,9,10,11,16,21,22,23,45$, but makes no claim for these four cases.
\end{enumerate}

Two further items concern the manuscript rather than the formalization.

\begin{enumerate}
\item Whether the equations and probability data developed here should be folded
  back into the manuscript. The arguments obtained for orbits 26, 29, 31 and 33
  are shorter than the manuscript's and in two cases avoid machinery it uses
  (matrix inverses and a genericity argument for orbit 26; a sixteen-type support
  classification for orbit 31), so there is a case for replacing rather than
  supplementing.
\item In a sibling companion manuscript, \path{../A_dyadic_ladder/paper/main.tex}: the
  sentence at line 1045 asserting that no lower bound is known is not correct for
  stabilizer codes, and the table caption at line 1000 needs revision. Neither
  touches the present development.
\end{enumerate}

\end{document}
